\scanpage{80}
\mspage{4}
% archive foliation: 79

est équivalente à celle des couples $(t,i)$, avec
\[
t\in U_{0,3}(k)=\mathbb{P}^1_k(k)\setminus\{0,1,\infty\},\quad\text{et }i\in\mu_4^*.
\]

\deleted{On peut définir sur $\Sigma_J$ le fibré}
\[
\deleted{\underline{O}_{\Sigma_J}(1)^{\unsure{\sharp}}\unsure{(}=\underline{O}_{\Sigma_J}(1)\wedge_{\unsure{\boldsymbol{\mu}|\mu}_4}}
\]
Je reviens à une base quelconque, mais je travaille en car. rés. $\ne2$ i.e. sur la base absolue\inserted{$S_0=$} $\operatorname{Spec}(\mathbb{Z}[\frac12])$. Sur \deleted{\unsure{$\Sigma_J$}} je définis, l'opération \deleted{\unsure{$\otimes$}} $\unsure{\sharp}$, sur les \deleted{\unsure{modules}} faisceaux en Modules \deleted{\unsure{sur}} sur un schéma de base $S$
\[
(22)\qquad M^{\unsure{\sharp}}=M\wedge_{\unsure{\boldsymbol{\mu}|\mu}_{4S}}\underbrace{Q_{0S}}_{\deleted{\unsure{\unsure{\boldsymbol{\mu}|\mu}}}}
\]
\margin{\uline{NB} $M^{\unsure{\sharp}}\simeq M\otimes_{\underline{O}_S}\underline{O}_S^{\unsure{\sharp}}$}
où ici on considère $Q_0$ comme un $\unsure{\boldsymbol{\mu}|\mu}_4$-torseur sur la base absolue $S_0$.

On aura donc
\[
(23)\qquad M^{\unsure{\sharp}}\otimes M^{\unsure{\sharp}}\simeq M^b\overset{\mathrm{def}}{=}M\wedge_{\unsure{\boldsymbol{\mu}|\mu}_{2S}=\{\pm1\}_S}\unsure{\boldsymbol{\mu}|\mu}_{4S}^*
\]
où
\[
\deleted{\unsure{\unsure{\boldsymbol{\mu}|\mu}_4^*\simeq}}
\]
\[
(24)\qquad\unsure{\boldsymbol{\mu}|\mu}_4^*=\unsure{\boldsymbol{\mu}|\mu}_4\setminus\unsure{\boldsymbol{\mu}|\mu}_2
\]
est considéré comme torseur sous $\unsure{\boldsymbol{\mu}|\mu}_2$, i.e. rev. étale quadratique de la base absolue $S_0$.

On s'intéresse donc, \deleted{\unsure{sur}} pour un revêtement
